\cleardoublepage{}

\thispagestyle{empty}
\null
\newpage

% ---------------------------------------------------------------------------
% Declaration of authorship (German, binding version) and English translation
% ---------------------------------------------------------------------------
\newcommand{\thesisauthor}{Alexander Giller}
\newcommand{\thesisplace}{München}
\newcommand{\thesisdate}{03.11.2026}

\thispagestyle{empty}
\vspace*{5mm}

\begin{center}
{\normalfont\bfseries Eidesstattliche Erklärung / Declaration of Authorship}
\end{center}

\vspace{8mm}

\noindent
\begin{otherlanguage}{ngerman}
Ich erkläre hiermit, dass ich die vorliegende Arbeit selbstständig und nur mit den angegebenen Quellen und Hilfsmitteln verfasst habe.
Alle Passagen, die ich wörtlich oder sinngemäß aus der Literatur oder aus anderen Quellen übernommen habe, sind als solche mit Quellenangabe gekennzeichnet.
Die Arbeit wurde in gleicher oder ähnlicher Form noch keiner anderen Prüfungsbehörde vorgelegt.
Diese gedruckte Fassung ist ein Ausdruck der eingereichten elektronischen Fassung.
\end{otherlanguage}

\vspace{6mm}

\noindent
I hereby declare that I have written this thesis independently and only with the sources and aids specified.
All passages that I have taken from literature or other sources, literally or in substance, have been marked as such with a reference to the source.
This thesis has not been submitted to any other examination authority in the same or a similar form.
This printed copy is a printout of the submitted electronic version.

\vspace{18mm}

\noindent
\begin{tabular}{@{}p{0.42\linewidth}@{\hspace{0.12\linewidth}}p{0.42\linewidth}@{}}
\thesisplace, \today & \\[-1ex]
\rule{\linewidth}{0.4pt} & \rule{\linewidth}{0.4pt} \\
Ort, Datum / Place, date & \thesisauthor
\end{tabular}

\vspace{14mm}

% ---------------------------------------------------------------------------
% Note on the use of AI tools (German and English)
% ---------------------------------------------------------------------------
{\small
\noindent
\begin{otherlanguage}{ngerman}
\textbf{Hinweis zur Nutzung von KI-Werkzeugen}

\noindent
Bei der Erstellung dieser Arbeit wurden Werkzeuge der künstlichen Intelligenz (KI), darunter große Sprachmodelle (LLMs) wie Claude Opus 5.5, zur Unterstützung bei Übersetzung und sprachlicher Überarbeitung eingesetzt, insbesondere bei Grammatik, Rechtschreibung, Klarheit und Lesbarkeit sowie bei der Umstrukturierung und Kürzung des Textes.
Darüber hinaus wurden KI-gestützte Werkzeuge für die Code-Vervollständigung, die Datenauswertung und die Erstellung von Abbildungen und Visualisierungen verwendet.
Alle KI-generierten oder KI-unterstützten Ergebnisse wurden vom Verfasser geprüft und kritisch bewertet. Der Verfasser trägt die volle Verantwortung für die Richtigkeit, die Interpretation und die endgültige Darstellung der Inhalte dieser Arbeit.
\end{otherlanguage}

\vspace{5mm}

\noindent
\textbf{Note on the Use of AI Tools}

\noindent
Artificial intelligence (AI) tools, including large language models (LLMs) such as Claude Opus 5.5, were used in the preparation of this thesis to support translation and linguistic editing, including improvements to grammar, spelling, clarity, and readability, as well as restructuring and condensing the text.
In addition, AI-based tools were used for code completion, data analysis, and the generation of figures and visualizations.
All AI-generated or AI-assisted outputs were reviewed and critically evaluated by the author. The author takes full responsibility for the accuracy, interpretation, and final presentation of the content of this thesis.
\par}

\cleardoublepage{}

\thispagestyle{empty}
\null
\newpage
